\fichatitulo{35 · ALINHAMENTO}{NeurIPS · 2022}{Training language models to follow instructions}
{\small Ouyang et al. \textit{Training Language Models to Follow Instructions with Human Feedback}. NeurIPS · 2022. \href{https://proceedings.neurips.cc/paper_files/paper/2022/file/b1efde53be364a73914f58805a001731-Paper-Conference.pdf}{Fonte consultada}.\par}

\campo{Problema e objetivo}
Como aproximar respostas das instruções e preferências humanas?

\campo{Principal contribuição · síntese interpretativa}
Combina demonstrações supervisionadas, rankings humanos e aprendizagem por reforço para ajustar modelos de linguagem.

\campo{Método / natureza da evidência}
Treina InstructGPT a partir de demonstrações e preferências; avalia respostas na distribuição de prompts utilizada no estudo.

\campo{Resultado ou argumento principal}
O modelo de 1,3 bilhão de parâmetros é preferido ao GPT-3 de 175 bilhões nessa distribuição. Preferência constitui o resultado específico comparado.

\campo{Excertos-chave · seleção literal}
\begin{itemize}
\excerto{1}{on our prompt distribution}{Resumo.}
\excerto{2}{InstructGPT still makes simple mistakes}{Resumo.}
\end{itemize}

\campo{Limitações e análise crítica}
Autores: persistem erros simples. Análise da ficha: preferência humana não equivale a verdade factual nem garante aderência a evidências parlamentares.

\campo{Aproveitamento na dissertação · proposta}
Fundamentação: diferenciar alinhamento, utilidade percebida e fidelidade às fontes recuperadas.

\registro{Fonte consultada nesta preparação: Resumo e introdução. Chave: \texttt{ouyangEtAl2022instructGpt}.}
